\documentclass[11pt,a4paper]{article}
\usepackage[margin=2.5cm]{geometry}
\usepackage[T1]{fontenc}
\usepackage[utf8]{inputenc}
\usepackage{lmodern}
\usepackage{setspace}
\usepackage{hyperref}
\usepackage{graphicx}
\usepackage{booktabs}
\usepackage{threeparttable}
\usepackage{amsmath, amssymb}
\usepackage{siunitx}
\usepackage{csquotes}
\usepackage{natbib}
\usepackage{xcolor}
\usepackage{enumitem}
\usepackage{caption}
\usepackage{subcaption}
\usepackage{placeins}
\usepackage{listings}
\lstset{basicstyle=\small\ttfamily,breaklines=true,columns=fullflexible}

\onehalfspacing
\setlength{\parskip}{0.4em}
\setlength{\parindent}{0pt}
\hypersetup{colorlinks=true,linkcolor=blue,citecolor=blue,urlcolor=blue}
\captionsetup{font=small}
\sisetup{detect-all=true}

\newcommand{\Wtwl}{\text{W12}}
\newcommand{\Wthfo}{\text{W34}}

\begin{document}

\begin{center}
{\LARGE Google Search Indicators and Nowcasting German Unemployment: Early vs. Late Information}\\[0.5em]
{\large Alex \\ Data Science in Business \& Economics}\\[0.25em]
{\small Supervisor: TBD \quad Seminar: TBD \quad Date: \today}
\end{center}
\vspace{1em}

\begin{abstract}
This paper evaluates whether Google search indicators aligned to institutional reporting windows improve one-step-ahead nowcasts of German unemployment relative to autoregressive benchmarks. Weekly search series (Indeed, Stepstone, Jobbörse, Arbeitsamt, Bewerbung) are stitched to consistent 2011--2023 histories and aggregated to bi-weekly components reflecting \Wthfo{} (M−1) and \Wtwl{} (M). We compare early vs. late nowcasts and assess gains via out-of-sample RMSE/MAE and Diebold–Mariano tests.
\end{abstract}

\section{Introduction}
The study examines the predictive content of high-frequency Google search data for nowcasting German unemployment. We introduce biweekly alignment reflecting institutional publication lags and contrast early versus late information windows.

\section{Related Literature}
Brief overview of literature on Google Trends for macroeconomic nowcasting and labor market indicators.

\section{Data}
\subsection{Target and Sample}
Unemployment rate (Bundesagentur für Arbeit) converted to monthly frequency for 2011--2023.
\begin{itemize}
    \item Problem of : The
announcement dates are provided in advance for the next two years and are almost always at the
end of the month, but sometimes early in the first week of the following month (Zimmerman)
\end{itemize}

\subsection{Google Search Indicators}
Weekly Google Trends indices for Indeed, Stepstone, Jobbörse, Arbeitsamt, and Bewerbung.
\textbf{Mention here}
\begin{itemize}
    \item Intro into google since when the engine exists 
    \item how you get info : A Google Insights query may have regional, temporal or keyword specific focus,
i.e. you choose the region of interest, the time frame of interest and the keywords of interest (up
to 5). The results are then delivered scaled and normalized within the query (for the region, the
time frame and the selection of keywords)2. 
\item Explaining the words and why I use them 
\item  What Zimmerman uses as timeframe and why
\end{itemize}
\subsection{Stitching Overlapping Windows}
Weekly windows are stitched via median-ratio chaining (W1→W2→W3) and aligned to a Sunday grid.
Check afterwards the distribution of the google trends.
Maybe as well correlations. Depending on what you get.

\subsection{Aggregation to Bi-Weekly Signals}
We create early- and late-month aggregates:
\begin{itemize}[leftmargin=*]
  \item $\textbf{W12(M)}$: mean of weekly Sundays with day $\le 15$.
  \item $\textbf{W34(M−1)}$ shifted to $M$: mean of Sundays with day $>15$ in $M-1$, shifted one month forward.
\end{itemize}
All series are month-end labeled and merged with unemployment.

\section{Methodology}
\subsection{Pre-tests and Transformations}
ADF and Engle–Granger tests are applied to each indicator and unemployment. Series without cointegration are modeled in first differences.

\subsection{Error--Correction Model (ECM)}
We estimate two variants:
\begin{align}
\Delta y_t &= \alpha + \phi\,\Delta y_{t-1} + \beta_0\,\Delta x^{\Wthfo}_t + \beta_1\,\Delta x^{\Wthfo}_{t-1} + \gamma\,EC^{early}_{t-1} + D_m + \varepsilon_t, \\
\Delta y_t &= \alpha + \phi\,\Delta y_{t-1} + \beta_0\,\Delta x^{\Wthfo}_t + \beta_1\,\Delta x^{\Wthfo}_{t-1} + \theta_0\,\Delta x^{\Wtwl}_t + \theta_1\,\Delta x^{\Wtwl}_{t-1} + \gamma\,EC^{late}_{t-1} + D_m + \varepsilon_t.
\end{align}

\subsection{Benchmarks and Evaluation}
AR and ARX baselines with month dummies; expanding-window OOS evaluation (36-month burn-in). Metrics: RMSE, MAE, DM tests, and encompassing checks.

\section{Results}
\subsection{Descriptive Checks}
Coverage and seasonality sanity checks; weekly grid continuity verified.

\subsection{Pre-test Outcomes}
Engle–Granger results indicate cointegration for Arbeitsamt\_W34prev and Jobbörse\_W34prev.

\subsection{Nowcast Performance}
Main results: in-sample RMSE/MAE diagnostics, encompassing regressions, and OOS accuracy with DM tests.

\subsection{Discussion}
Interpretation of early vs. late results and adjustment speeds; institutional timing and robustness.

\section{Robustness}
\subsection{Pairwise Engle--Granger Screen}
Pairwise Engle--Granger tests using different trends ($c$, $ct$) confirm strongest cointegration for Arbeitsamt\_W34prev and Jobbörse\_W34prev. Others remain I(1) and serve as short-run drivers.

\subsection{Additional Checks}
Alternative lags, month-dummy controls, PCA combinations, and sample-split analyses (COVID subperiod).

\section{Conclusion}
Google search indicators provide timely signals of labor market dynamics. Late-month information marginally improves nowcasts when combined in ECM frameworks.

\section*{References}
\bibliographystyle{apalike}
\bibliography{seminar}

\appendix
\section{Data Construction Details}
Pseudocode and code excerpts for stitching and aggregation.

\section{Additional Tables and Figures}
Additional diagnostics, residual checks, and robustness tables.

\subsection{Dynamic Interpretation and Extensions}
% (This subsection expands ideas (a)–(e) for final text)

\paragraph{(a) Economic interpretation of coefficients.}
Interpret the sign and size of $\gamma$ (error–correction speed): how fast unemployment returns to equilibrium.
Discuss $\beta_0,\beta_1$ (short-run impact of late-month W34prev activity)
versus $\theta_0,\theta_1$ (impact of early-month W12 searches).

\paragraph{(b) Dynamic simulation.}
Use estimated ECM parameters to simulate a shock scenario (e.g., a 10\% rise in Stepstone W34prev).
Plot the cumulative response of $\Delta \text{Unemp}$ and the implied level path over three months.

\paragraph{(c) Cross–indicator comparison.}
Contrast results across \textit{Arbeitsamt}, \textit{Jobbörse}, \textit{Indeed}, and \textit{Stepstone}.
Highlight which represent public inflow versus private outflow searches and how that affects predictive power.

\paragraph{(d) Seasonality and policy timing.}
Summarize and interpret month–dummy coefficients.
Relate recurring seasonal dips (e.g., February–September) to hiring cycles, holidays, and institutional effects.

\paragraph{(e) Model comparison and parsimony.}
Compare AIC/BIC between Early and Late ECMs.
Discuss potential overfitting when adding W12 terms, and justify the preferred specification.

\end{document}